\begin{conclusion}
本课题围绕资源受限场景下 VL-MoE 模型的推理优化，完成了一个面向 Qwen3-VL 的可复现实验原型。整体工作并不只停留在推理瓶颈分析，而是从真实链路 profiling、视觉侧 Token 压缩、异步流水执行、专家路由 trace 采集、全局预测器训练，到 CPU-GPU expert offload 调度验证，形成了较完整的系统实现流程。

从研究定位看，本文关注的不是单一算子的局部加速，而是多模态稀疏模型在真实推理链路中的复合瓶颈。VL-MoE 同时包含视觉前端、语言主干、稀疏专家和动态调度等环节，任何一个环节的等待都可能放大端到端 TTFT。因此，本文将视觉 Token 冗余和专家权重访问动态性作为两个主要问题，并分别设计语义感知压缩和预测式专家调度进行处理。

在视觉侧，本课题实现了语义感知视觉 Token 压缩模块。该模块利用文本语义对视觉区域进行相关性引导，并结合自适应压缩率控制视觉 Token 数量，从而减少 Visual Encoder 后续层的冗余计算。在此基础上，系统进一步实现了视觉侧与文本侧准备过程的异步流水执行，通过双流并发隐藏部分视觉计算等待。实验结果表明，在 MMBench 与 MME 子集上，优化后的归一化任务分数均保持在 0.995 以上，说明当前压缩策略没有造成明显效果下降；同时，Visual Encoder 可见耗时由 8572ms 降至 800ms，整体 TTFT 获得约 1.63x 加速。

在专家侧，本课题完成了 Qwen3-VL fresh 路由 trace 的采集与训练数据构建，并实现了 layer-0 条件化的全局路由预测器。预测器输入包括首层 router 概率、Top-$k$ 专家、token 位置、输入 id、模态标识和路由统计特征，用于预测后续层更可能访问的专家集合。训练结果显示，预测器在验证集上的 best mean recall 达到 0.7219，说明 VL-MoE 的专家路由具有可利用的跨层结构性。

围绕预测结果，本课题还实现了预测驱动的受限专家预取与缓存调度机制。系统根据候选专家置信度、层距衰减、预测 margin 和缓存预算选择预取对象，并通过保护当前访问专家和同步回退机制保证推理过程能够持续执行。在真实 CPU-GPU expert offload 微基准中，predictor 策略在 \texttt{cap96} 和 \texttt{cap160} 两个代表性设置下分别取得 2.99x 和 2.66x 加速，验证了提前调度专家权重能够显著降低同步加载等待。

综合来看，本文的主要贡献可以概括为三点。第一，基于真实模型和公开数据集对 VL-MoE 推理链路进行了较系统的瓶颈分析，指出视觉侧前置开销和专家侧动态访问共同限制推理效率。第二，提出并实现了面向 Visual Encoder 的语义感知 Token 压缩与异步流水方法，在效果保持和时延下降之间取得了较好的平衡。第三，构建了面向 Qwen3-VL 的 fresh 路由预测数据集，并验证了 layer-0 条件化预测在专家预取和缓存调度中的实用价值。

目前系统仍存在一些不足。首先，视觉侧优化与专家侧调度虽然分别完成了实现和实验验证，但二者尚未在稳定的完整端到端 serving 框架中充分融合；现有基于 Transformers generate 的端到端原型可以运行，但受到 Python hook、图像预处理和缓存控制等因素影响，结果波动较大。其次，路由预测数据规模仍然有限，主要覆盖当前模型和实验数据集，预测器在更多任务类型、更长上下文和不同 VL-MoE 架构上的泛化能力还需要进一步验证。再次，专家预取调度仍带有原型系统特征，部分控制逻辑位于 Python 层，缓存替换、传输重叠和预算自适应还不够精细；在缓存容量过紧或预测置信度不足时，额外预取可能带来无效传输，抵消部分加速收益。

后续工作可以从三个方向继续完善。第一，扩大路由 trace 的采集规模，引入更多任务类型、模态比例和上下文长度，并尝试更强的时序建模方法，以提升全局专家预测的稳定性和泛化能力。第二，将专家预取、缓存替换和异步传输进一步下沉到更低开销的 runtime 层，减少 Python 侧调度开销，并实现按请求动态调整的缓存预算和预取阈值。第三，构建更完整的端到端部署评估环境，将视觉 Token 压缩、异步流水、专家预测调度与批处理、KV cache 管理等服务端机制结合起来，验证系统在真实多用户负载下的延迟、吞吐和显存收益。

除此之外，视觉侧优化也可以继续向更细粒度方向发展。当前 SAT-ToMe 主要根据全局文本语义和视觉 Token 相似度进行选择，对于多图像输入、视频帧输入和高分辨率文档页面，未来可以引入区域级预算、跨图像均衡保留和层级压缩机制，避免少数显著区域占用过多 Token 预算。专家侧则可以进一步研究模态分组预测、在线反馈校准和多请求共享预取，使缓存策略能够适应在线服务中不断变化的请求分布。

总体来看，本课题完成了 VL-MoE 推理优化原型中从瓶颈定位到关键模块实现、再到真实 offload 微基准验证的主要工作。当前系统已经证明语义感知视觉压缩和预测式专家调度具有实际加速潜力，而后续需要重点补足端到端集成、跨任务泛化和 runtime 工程化能力，使其从实验原型进一步走向更稳定的部署系统。
\end{conclusion}
